% SOURCE_STATEMENT: "A number of so-called ``transportation'' problems have now been analyzed in terms of RMT. These include: the ``vicious'' walker problem of M. Fisher, the bus problem in Cuernavaca, Mexico, the headway traffic problem on highways, and the airline boarding problem of Bachmat et al. Recently, researchers in London, Prague, and also Ann Arbor, have noticed an intriguing phenomenon. They found that the fluctuations in the spacings between cars parked on a long street exhibited RMT behavior. Furthermore, \\v{S}eba found that there was a difference whether the street is two-way or one-way (On a two-way street, the cars park only on the right, while on a one-way street one of course has the option of also parking on the left.) Quite remarkably, for two-way streets \\v{S}eba found GUE statistics, but for left-side parking on one-way streets he found GOE statistics. It is a great challenge to develop a microscopic model for the parking problem, in analogy, perhaps, with the microscopic model introduced by Baik et al. to explain the RMT statistics for the bus problem in Cuernavaca. \\v Seba's recent, intriguing calculations on the parking problem can be found posted on the web."
% Preparation and final self-audit: CONTRIBUTING.md prompts C and D; source/audit artifacts retained in the local physics_statements_2026_10_11 batch.
% !TeX program = lualatex
\documentclass[11pt]{article}
\usepackage[margin=25mm]{geometry}
\usepackage{fontspec}
\setmainfont{Latin Modern Roman}
\usepackage{amsmath,amssymb,mathtools}
\usepackage{unicode-math}
\setmathfont{Latin Modern Math}
\usepackage{xurl,hyperref}
\hypersetup{colorlinks=true,urlcolor=blue}
\setcounter{secnumdepth}{0}
\setlength{\parindent}{0pt}
\setlength{\parskip}{5pt}
\setlength{\emergencystretch}{3em}
\newcommand{\sref}[2]{\hyperlink{source:#1:#2}{[S#2]}}
\newcommand{\cataloguescope}[1]{\paragraph{Recorded target.}#1}
\newcommand{\R}{\mathbb R}
\newcommand{\C}{\mathbb C}
\newcommand{\N}{\mathbb N}
\newcommand{\Z}{\mathbb Z}
\newcommand{\Q}{\mathbb Q}
\newcommand{\F}{\mathbb F}
\newcommand{\eps}{\varepsilon}
\begin{document}
% UnsolvedMath ID 8000009; source catalogue code AMR-079-0009
\setcounter{section}{8000008}
\section{The parking problem}\label{8000009}
{\small UnsolvedMath ID 8000009 · Mathematical Physics}
\subsection{Definitions and mathematical statement}

\paragraph{Shared notation.}$\mathbb N=\{1,2,\ldots\}$; the standard number fields have their usual meanings. All additional conventions are stated locally.


\paragraph{Definitions and notation.}The target concerns observed spacings between real cars parked on a long street. A microscopic stochastic model specifies the actions and interactions of individual drivers and cars, from which spacing statistics are derived. GUE and GOE are Gaussian unitary and orthogonal random-matrix ensembles; their normalized local-spacing laws are different. Random sequential interval deposition is one possible parking model, not the model prescribed by this source.

\paragraph{Catalogue wording retained for traceability.}
Develop a microscopic model for parked-car spacings explaining the source’s observed random-matrix behavior: GUE statistics for two-way-street parking on the right and GOE statistics for left-side parking on one-way streets.
\paragraph{Scope and source qualification.}The source refers to empirical spacing studies by Šeba and others. A theorem for the classical random sequential adsorption model alone would be a different target unless it explains those observations.
\subsection{Short English statement}
Build a microscopic model explaining the different random-matrix spacing statistics observed for parked cars.
\subsection{Sources}
\begin{itemize}
\item[\hypertarget{source:8000009:1}{S1}] ULAM.AI and the UnsolvedMath Contributors, UnsolvedMath, record 8000009 (AMR-079-0009). Catalogue curation under CC BY 4.0; definitions and English exposition added locally.
\url{https://huggingface.co/datasets/ulamai/UnsolvedMath}.
\item[\hypertarget{source:8000009:2}{S2}] Percy Deift, Some open problems in random matrix theory and the theory of integrable systems (2007), arXiv:0712.0849, Problem 9.
\url{https://arxiv.org/abs/0712.0849}.
\end{itemize}
\label{end:8000009}
\end{document}
